% Standalone section -- paste into the weekly report, or \input it.
% Week 3 working note: whether local branch volume at a bifurcation separates healthy from
% diseased cases, and a candidate mechanism (Glagov remodeling) for why the sign of that
% separation flips between bifurcations. Companion to bifurcation_angle_and_disease.tex, which asks
% the same disease-separation question of angle instead and comes back null. Numbers from
% scripts/bifurcation_volume.py, run 2026-09-16 (800 cases); assumes the tree and angle/caliber
% definitions from coronary_tree_graph.tex and bifurcation_angle.tex.

\section*{Local branch volume separates healthy from diseased cases}

Bifurcation angle, measured in bifurcation\_angle.tex, is a direction: it says nothing about how
much lumen the two daughter branches actually carry. The caliber panel already read alongside that
angle, radius ratio and the two branches' absolute radii, hints that size at the junction is where
the disease signal might actually sit. This note turns that hint into a proper test by extending
the same measurement past a single cross-section into a volume: \textbf{local branch volume}, the
frustum-integrated lumen volume of the main and side daughter over the same shaft window the angle
and radius are already read from (past the bifurcation core, out to 3\,mm of arc). Reusing that
window rather than defining a new one keeps volume directly comparable to the angle and caliber
numbers already in hand, case for case and bifurcation for bifurcation, rather than a whole-vessel
volume that would mix in geometry the angle measurement never sees.

\subsection*{Setup: integrating a frustum along the shaft}

\texttt{scripts/bifurcation\_volume.py} walks the same shaft window \texttt{topology/angles.py}
uses for direction and caliber and, between each pair of consecutive centerline points in that
window, treats the vessel as a frustum (a cone with the top cut off) of the two points' radii and
sums $\frac{\pi}{3} d (r_0^2 + r_0 r_1 + r_1^2)$ over every step, $d$ being the distance between the
points. That is exact for a straight-walled cone and a close approximation everywhere the radius
changes smoothly point to point, which the shaft window is chosen to guarantee by construction. The
same window is walked for the main branch and the side branch separately, and their sum, on a log
scale, is also tested as a single per-bifurcation size variable (\texttt{v\_total}).

Every comparison below is Mann--Whitney $U$ on the per-case mean of that quantity, healthy against
diseased, with Benjamini--Hochberg FDR correction applied across all 18 bifurcation$\times$field
combinations run together, not per test, so that finding six real effects among eighteen is not
easier by chance than it looks.

\subsection*{Result: three of five bifurcations separate, after correction}

Local branch volume is larger in diseased cases at two of five named bifurcations and smaller at a
third, and all three survive that correction (\Cref{tab:week3-volume}). Neither the
daughter-daughter angle nor its parent-relative split (\texttt{theta\_main\_deg},
\texttt{theta\_side\_deg}) survives correction at any of the 48 angle/caliber tests run the same
way, at any bifurcation, so the separation volume finds is not simply size finding what angle also
finds; the two quantities disagree about which cases are diseased. The companion note,
bifurcation\_angle\_and\_disease.tex, works through the angle side of that disagreement in full,
including whether adding volume to angle recovers a signal angle lacks on its own (it does not: at
CRUX, angle alone predicts disease at chance while volume alone reaches AUC 0.576, and angle added
to volume does not improve on volume).

\begin{table}[h]
  \centering
  \caption{Diseased vs.\ healthy local main-daughter branch volume, per-case mean, Mann--Whitney
    $U$ with Benjamini--Hochberg FDR across all 18 bifurcation\,$\times$\,volume-field tests run.
    Only the six surviving $q < 0.05$ are shown; the pooled-over-all-bifurcations and LM/LAD-D2
    tests did not (\texttt{scripts/bifurcation\_volume.py}).}
  \label{tab:week3-volume}
  \begin{tabular}{lrrrrr}
    \hline
    Bifurcation & $n$ healthy & $n$ diseased & healthy (mm$^3$) & diseased (mm$^3$) & $q$ \\
    \hline
    CRUX (PDA/PLA) & 400 & 371 & 6.74 & 7.78 (+15.5\%) & 0.003 \\
    LCX/OM1        & 352 & 342 & 7.59 & 8.45 (+11.3\%) & 0.036 \\
    LAD/D1         & 407 & 380 & 10.80 & 10.02 ($-$7.2\%) & 0.036 \\
    \hline
  \end{tabular}
\end{table}

The disease-status label used throughout (\texttt{Disease}, alongside \texttt{Image Quality} and
\texttt{Dominance}) is ImageCAS-X's own binary annotation, present or absent, with no severity
grade and no record of which vessel carries it \cite{bransby2026imagecasx}; every comparison in
this note inherits that granularity, and the discussion below returns to what it costs.

\subsection*{The signal is not an artifact of mixing dominance}

The crux is measured at a different pair of vessels depending on dominance, R-PDA/R-PLA or
L-PDA/L-PLA, so a pooled CRUX comparison risks averaging over two anatomically distinct junctions,
and the worry going in was that this alone might be manufacturing the effect. Splitting by the
cohort's own dominance annotation resolves this rather than weakening the result: restricted to the
721 right-dominant cases, a single fixed anatomical site in every one of them, the CRUX effect is
stronger, not weaker ($+17.1\%$, $p = 6.1\times10^{-5}$, $n_h=374$, $n_d=347$), and LCX/OM1
replicates the same direction ($+12.0\%$, $p = 0.0089$, $n_h=316$, $n_d=307$). The 32 left-dominant
and 30 codominant cases split too thinly by bifurcation ($n \le 19$ per disease group) to test on
their own, and what little signal is there trends the opposite way; whether the effect is
dominance-independent and simply underpowered in those two strata, or a real interaction with
dominance, is open and cannot be settled inside ImageCAS-X.

\subsection*{A candidate mechanism: compensatory remodeling, not a fixed reading}

The sign flip in \Cref{tab:week3-volume}, larger at two bifurcations and smaller at a third, needed
some explanation before it could be trusted rather than shrugged off as noise, and the most direct
one on offer is arterial remodeling. Glagov et al.\ examined histologic sections of the left main
from 136 hearts at autopsy and found the artery wall enlarges as plaque accumulates: wall area
correlated with lesion area ($r = 0.44$, $p < 0.001$), and lumen area did not shrink while the
lesion occupied up to 40\% of the internal elastic lamina area, falling sharply only above that
threshold ($r = -0.73$, $p < 0.001$) \cite{glagov1987remodeling}. This licenses a specific reading
of the sign flip: CRUX and LCX/OM1, where diseased volume is larger, sit in the compensated regime;
LAD/D1, where it is smaller, has exceeded it, or has less remodeling reserve to begin with. It does
not license treating this as settled: ImageCAS-X's binary \texttt{Disease} label carries no
plaque-burden fraction, so the 40\% threshold itself cannot be tested against our own cases, only
proposed as the mechanism a burden variable would need to check. Glagov's finding also states the
caveat that governs every comparison in this note: a lumen segmentation has already been reshaped
by the disease it is being asked to predict, so a cross-sectional geometry--disease association is
partly the disease writing its own predictor rather than a property that preceded it.

\subsection*{What this licenses}

Local branch volume, not angle, is the ImageCAS-X-level disease signal worth carrying forward, with
Glagov's compensatory-remodeling law as a stated, not yet tested, mechanism for why it changes sign
between bifurcations. What it does not license is reading that sign flip as causal or as a stage
marker: without a plaque-burden fraction to place each case on Glagov's curve, the same 800 cases
that produced \Cref{tab:week3-volume} cannot distinguish "compensated" from "decompensated" any
other way. The ramus intermedius natural experiment (\texttt{IM} label, Section~\ref{sec:tree}) and
the CGPS baseline-vs-progression design already in the project plan are where a severity-graded or
longitudinal test of this mechanism becomes possible.
